\begin{theorem}[complete Cartier-fibre duality]
\hypertarget{tag-03IY}{}\label{wbII:thm:cartier-duality}
\StatementRecord{wbII:thm:cartier-duality}{proved}{03IY}{\cite[Proposition~2 and pp.~90--91]{RR70push}; \cite[Theorem~2, pp.~82--83]{RR70push}}

\ Let $X$ be a compact smooth complex threefold,
$i:W\hookrightarrow X$ an effective Cartier divisor, and $E$ a holomorphic
vector bundle.  Put $F=E|_W$ and
\[
\omega_W=(\omega_X\otimes\mathcal O_X(W))|_W.
\]
For $q=0,1,2$ there are natural perfect pairings
\[
H^q(W,F)^\vee
\simeq H^{2-q}(W,F^\vee\otimes\omega_W).
\tag{W.11}\hypertarget{tag-03EB}{}\label{wbII:eq:W:5.11}
\]
Moreover $H^q(W,F)=0$ for $q>2$.  No normality or reducedness hypothesis on
$W$ is used.
\end{theorem}



\begin{proof}

\begin{enumerate}[label={$\langle1\rangle$\arabic*.},leftmargin=*]
\item Let \(i:W\hookrightarrow X\) be the given Cartier inclusion.  The
length-one locally free resolution
\[
 0\longrightarrow\mathcal O_X(-W)
 \xrightarrow{\;\cdot g\;}\mathcal O_X
 \longrightarrow i_*\mathcal O_W\longrightarrow0
\]
is locally defined by a nonzerodivisor \(g\); changing \(g\) by a unit
gives the usual transition functions for \(\mathcal O_X(W)\).
Applying sheaf Hom into \(\omega_X\), the degree-zero map is
injective and its cokernel is
\[
 \mathcal E xt^1_X(i_*\mathcal O_W,\omega_X)
 =i_*\bigl((\omega_X\otimes\mathcal O_X(W))|_W\bigr)
 =i_*\omega_W.
\]
The other sheaf Ext groups vanish.  Equivalently, \(W\) is a
two-dimensional Cohen--Macaulay Gorenstein complex space and its
analytic dualizing complex is
\[
 K_W^\bullet\simeq\omega_W[2].
\]
This is the hypersurface calculation in
\cite[Proposition~2 and pp.~90--91]{RR70push}; no normality or
reducedness of the Cartier divisor is used.

\item For the locally free sheaf \(F=E|_W\), the analytic duality trace of
\cite[Theorem~2, pp.~82--83]{RR70push} supplies the actual
functorial pairing, and identifies the dual of \(H_c^q(W,F)\) with
\[
 \operatorname{Ext}^{-q}_W(F,K_W^\bullet)
 =\operatorname{Ext}^{2-q}_W(F,\omega_W)
 =H^{2-q}(W,F^\vee\otimes\omega_W).
\]
The last equality follows because the positive sheaf Ext groups
of a locally free first argument vanish.  Since \(W\) is compact,
\(H_c^q(W,F)=H^q(W,F)\); coherent cohomology is finite-dimensional,
so this pairing is perfect in the ordinary algebraic sense.
It is this trace pairing, rather than an unspecified arrow between
two matching five-term sequences, that constructs the isomorphism
\eqref{wbII:eq:W:5.11}.  Its contravariance in \(F\) is precomposition in the Ext
argument.

\item For \(q>2\), the displayed Ext degree \(2-q\) is negative.  The
derived Hom of the locally free sheaf \(F\), concentrated in degree
zero, into the sheaf \(\omega_W\), concentrated in degree zero, has
no negative cohomology.  Hence the duality theorem gives
\(H^q(W,F)=0\).  This also follows directly from coherent
cohomological dimension on the compact surface \(W\).

\end{enumerate}
\end{proof}